%───────────────────────────────────────────────────────────────────────────────
% TABLES-CONFIG.TEX — Estilo global de tablas con paleta Catppuccin
%───────────────────────────────────────────────────────────────────────────────
% Requiere: booktabs, array, tabularx, longtable, multirow, siunitx,
%           colortbl, xparse, etoolbox
% colortbl se carga en packages.tex
%───────────────────────────────────────────────────────────────────────────────

% ─── Espaciado global de booktabs ────────────────────────────────────────────
\setlength{\heavyrulewidth}{0.08em}
\setlength{\lightrulewidth}{0.04em}
\setlength{\cmidrulewidth}{0.03em}
\renewcommand{\arraystretch}{1.25}

% ─── Alineación vertical: m (middle) en lugar de p (top) ────────────────────
% Por defecto, tabularx usa p{} (alineación al tope) para las columnas X.
% Esto causa desalineación vertical cuando se mezclan columnas l/c/r con X.
% Redefinimos para que X use m{} (centrado vertical):
\renewcommand{\tabularxcolumn}[1]{m{#1}}

% ─── Nuevos tipos de columna ─────────────────────────────────────────────────
% L/C/R: columnas de ancho fijo con wrapping y centrado vertical
\newcolumntype{L}[1]{>{\raggedright\arraybackslash}m{#1}}
\newcolumntype{C}[1]{>{\centering\arraybackslash}m{#1}}
\newcolumntype{R}[1]{>{\raggedleft\arraybackslash}m{#1}}

% Y/Z: columnas proporcionales para tabularx (centrada / derecha)
% X ya existe en tabularx (izquierda, proporcional, ahora con m{})
\newcolumntype{Y}{>{\centering\arraybackslash}X}
\newcolumntype{Z}{>{\raggedleft\arraybackslash}X}

% ─── Colores para tablas ─────────────────────────────────────────────────────
\colorlet{tableheadcolor}{CtpBlue!12}
\colorlet{tableheadtext}{CtpText}
\colorlet{tablerowalt}{CtpSurface0!30}
\colorlet{tablerule}{CtpOverlay0!50}

% Aplicar color a reglas de booktabs (después de \begin{document})
\AtBeginDocument{\arrayrulecolor{tablerule}}

% ─── Comando: encabezado de columna ─────────────────────────────────────────
\newcommand{\thead}[1]{%
  \bfseries\color{tableheadtext}#1%
}

% Comando para colorear la fila de encabezado completa
\newcommand{\headerrow}{\rowcolor{tableheadcolor}}

% ─── Entorno: tabla con overflow controlado ──────────────────────────────────
% Reescala tablas anchas para que quepan en \textwidth.
% Uso:
%   \begin{fittable}
%     \begin{tabular}{llllll}  ...  \end{tabular}
%   \end{fittable}
%
\newsavebox{\fittablebox}
\NewDocumentEnvironment{fittable}{}{%
  \begin{lrbox}{\fittablebox}%
}{%
  \end{lrbox}%
  \ifdim\wd\fittablebox>\textwidth
    \resizebox{\textwidth}{!}{\usebox{\fittablebox}}%
  \else
    \usebox{\fittablebox}%
  \fi
}

% ─── Entorno: tabla flotante con caption ─────────────────────────────────────
% Uso:
%   \begin{acadtable}[htbp]{tab:label}{Caption text}
%     \begin{tabularx}{\textwidth}{lYY}  ...  \end{tabularx}
%   \end{acadtable}
%
\NewDocumentEnvironment{acadtable}{ O{htbp} m m }{%
  \begin{table}[#1]%
    \centering
    \caption{#3}\label{#2}%
}{%
  \end{table}%
}

% ─── Caption: renombrar "Cuadro" → "Tabla" en español ───────────────────────
\addto\captionsspanish{\renewcommand{\tablename}{Tabla}}

% ─── Notas al pie de tabla ───────────────────────────────────────────────────
\newcommand{\tablenote}[2]{%
  \par\vspace{2pt}%
  {\footnotesize\textsuperscript{#1}\,#2}%
}

%───────────────────────────────────────────────────────────────────────────────
% NOTA: tabularx NO se puede envolver en \NewDocumentEnvironment porque
% internamente escanea el body buscando \end{tabularx} de forma literal.
% Usar directamente:
%   \begin{tabularx}{\textwidth}{l Y Y Z}
%     ...
%   \end{tabularx}
%
% Las columnas Y (centrada) y Z (derecha) se expanden proporcionalmente.
%───────────────────────────────────────────────────────────────────────────────
